\subsection{Connected components of a quiver}

Let \(C = (S, F, o, t)\) be a quiver. Consider the equivalence relation \(R_{C}\) on S, the finest such that the relation \(\mathrm{R}_{\mathrm{C}}\{o(f), t(f)\}\) is satisfied for every arrow f of C. Two vertices a, b of C are equivalent for this relation if and only if there exists an integer \(n \geqslant 0\) , vertices \(a_{0}, \ldots, a_{n}\) of C and arrows \(f_{1}, \ldots, f_{n}\) of C such that \(a_{0} = a\) , \(a_{n} = b\) and such that, for \(1 \leqslant i \leqslant n\) , the arrow \(f_{i}\) joins either \(a_{i-1}\) to \(a_{i}\) , or \(a_{i}\) to \(a_{i-1}\) .

The equivalence classes of this equivalence relation are called the connected components of C. We denote by \(\pi_0(\mathrm{C})\) the set of connected components of C. Finally, we say that C is connected if it has at most one connected component.

Let \(\varphi\colon\mathrm{C}\to\mathrm{C}^{\prime}\) be a morphism of quivers. The map \(\operatorname{Som}(\varphi)\) defines, by passing to quotients, a map from \(\pi_{0}(\mathrm{C})\) to \(\pi_{0}(\mathrm{C}^{\prime})\) which we denote \(\pi_{0}(\varphi)\) .

